\documentclass[../../main.tex]{subfiles}
\begin{document}

\begin{center}
  {\fontsize{48pt}{58pt}\selectfont \textbf{1966}\normalsize\textbf{年度}}
\end{center}

\begin{table}[htb]
  \centering
  \caption{出題テーマと難易度}
  \begin{tabular}{|c|p{9cm}|c|}
    \hline
    \textbf{問題番号} & \textbf{テーマ} & \textbf{難易度} \\
    \hline
    第1問 & 三次方程式の実根と無理性 & \\
    \hline
    第2問 & 平行四辺形に内接する三角形の最大面積 & \\
    \hline
    第3問 & 新課程：塔の見え方と共線性／旧課程：4点間の距離 & \\
    \hline
    第4問 & 新課程：回転する単位ベクトル間の距離／旧課程：二次方程式の実根 & \\
    \hline
    第5問 & 新課程：指数変換された正方形の周の像／旧課程：関数の最大値の積分 & \\
    \hline
    第6問 & 新課程：正多角形の頂点から作る鋭角三角形／旧課程：正数列 & \\
    \hline
  \end{tabular}
\end{table}

\begin{table}[htb]
  \centering
  \caption{解答}
  \renewcommand{\arraystretch}{1.4}
  \begin{tabular}{|c|c|p{9cm}|}
    \hline
    \textbf{問題番号} & \textbf{小問} & \textbf{答え} \\
    \hline
    第1問 & (1) & $1<x<2$にただ1つ実根をもつ（証明問題） \\
    \cline{2-3}
    & (2) & その根は無理数（証明問題） \\
    \hline
    第2問 & (1) & 最大面積の三角形は周上に頂点があり，少なくとも一辺が平行四辺形の一辺と一致する \\
    \cline{2-3}
    & (2) & 面積$1$の三角形は面積$2$未満の平行四辺形に入らない \\
    \hline
    第3問 & 新課程 & $A,B,C$は一直線上にある（証明問題） \\
    \cline{2-3}
    & 旧課程 &  \\
    \hline
    第4問 & 新課程 & $0<a\le\dfrac{1}{5}$では$2-a$，$a\ge\dfrac{1}{5}$では$\sqrt{a^2+\dfrac{9a}{4}+\dfrac{3}{2}+\dfrac{1}{4a}}$ \\
    \cline{2-3}
    & 旧課程 &  \\
    \hline
    第5問 & 新課程 & $2(a^3-a^{-3})\log a$ \\
    \cline{2-3}
    & 旧課程 &  \\
    \hline
    第6問 & 新課程(1) & $\dfrac{n(n-1)(n-2)}{6}$ \\
    \cline{2-3}
    & 新課程(2) & $n$が奇数なら$\dfrac{n(n-1)(n+1)}{24}$，偶数なら$\dfrac{n(n-2)(n-4)}{24}$ \\
    \cline{2-3}
    & 旧課程 &  \\
    \hline
  \end{tabular}
\end{table}

\paragraph{試験形態}
6問．試験時間は要確認．

\end{document}
